\section{The 2-category of kernels}

In this section we introduce the 2-category $\gVar$, which provides
the natural context for the study of the structure of integral
transforms between derived categories of spaces.  The objects of
$\gVar$ are spaces, 1-morphisms are kernels of integral transforms,
and 2-morphisms are maps between these kernels.  Before introducing
$\gVar$ we remind the reader of the notion of a 2-category and we
explain the string diagram notation of which we will have much use.

 \subsection{A reminder on 2-categories.}
We will review the notion of a 2-category at the same time as
introducing the notation we will be using. Recall that a 2-category
consists of three levels of structure: objects; 1-morphisms between
objects; and 2-morphisms between 1-morphisms.  It is worth mentioning
a few examples to bear in mind during the following exposition.

\begin{itemize}
\item[1.] The first example is the 2-category $\gCat$ of categories,
  functors and natural transformations.
\item[2.] The second example is rather a family of examples.  There is
  a correspondence between 2-categories with one object $\star$ and
  monoidal categories.  For any monoidal category the objects and
  morphisms give respectively the 1-morphisms and 2-morphisms of the
  corresponding 2-category.
\item[3.] The third example is the 2-category $\gAlg$ with algebras
  over some fixed commutative ring as its objects, with the set of
  $A$-$B$-bimodules as its 1-morphisms from $A$ to $B$, where
  composition is given by tensoring over the intermediate algebra, and
  with bimodule morphisms as its 2-morphisms.
\end{itemize}


There are various ways of notating 2-categories: the most common way
is to use \emph{arrow diagrams}, however the most convenient way for
the ideas in this paper is via \emph{string diagrams} which are
Poincar\'e dual to the arrow diagrams.  In this subsection we will
draw arrow diagrams on the left and string diagrams on the right to
aid the reader in the use of string diagrams.


Recall the idea of a 2-category.  For any pair of objects $X$ and $Y$
there is a collection of morphisms $\OHom(X,Y)$; if
$\Phi\in\OHom(X,Y)$ is a 1-morphism then it is drawn as below.
\[\pic{AD1}\qquad\pic{SD1}\]
These 1-dimensional pictures will only appear as the source and target
of 2-morphism, i.e., the top and bottom of the 2-dimensional pictures
we will be using.  In general 1-morphisms will be denoted by their
identity 2-morphisms, see below.


If $\Phi,\Phi'\in\OHom(X,Y)$ are parallel 1-morphisms --- meaning
simply that they have the same source and target --- then there is a
set of 2-morphisms $\THom(\Phi,\Phi')$ from $\Phi$ to $\Phi'$.  If
$\alpha\in\THom(\Phi,\Phi')$ is a 2-morphism then it is drawn as
below.
\[\pic{AD2}\qquad\pic{SD2}\]
At this point make the very important observation that diagrams are read from right to
left and from bottom to top.


There is a \emph{vertical composition} of 2-morphisms so that if
$\alpha\colon\Phi\Rightarrow\Phi'$ and
$\alpha'\colon\Phi'\Rightarrow\Phi''$ are 2-morphisms then the
vertical composite $\alpha'\circv\alpha\colon\Phi\Rightarrow\Phi''$ is
defined and is denoted as below.
\[\pic{AD3a}\equiv\pic{AD3b}\qquad\pic{SD3}\equiv\pic{SD3b}\]
This vertical composition is strictly associative so that
$(\alpha''\circv\alpha')\circv\alpha
=\alpha''\circv(\alpha'\circv\alpha)$ whenever the three 2-morphisms
are composable.  Moreover, there is an identity 2-morphism
$\id_\Phi\colon \Phi\Rightarrow \Phi$ for every 1-morphism $\Phi$ so
that $\alpha\circv\id_\Phi=\alpha=\id_{\Phi'}\circv\alpha$ for every
2-morphism $\alpha\colon\Phi\Rightarrow\Phi'$.  This means that for
every pair of objects $X$ and $Y$, the 1-morphisms between them are the
objects of a category $\bHom(X,Y)$, with the 2-morphisms forming the
morphisms.  In the string notation the identity 2-morphisms are usually just
drawn as straight lines.
\[\pic{AD4}\qquad\pic{SD4}\]


There is also a composition for 1-morphisms, so if $\Phi\colon X\to Y$
and $\Psi \colon Y\to Z$ are 1-morphisms then the composite $\Psi\circ
\Phi\colon X\to Z$ is defined and is denoted as below.
\[\pic{AD5a}\equiv\pic{AD5b}\qquad\pic{SD5}\equiv\pic{SD5b}\]
Again, these pictures will only appear at the top and bottom of
2-morphisms.  This composition of 1-morphisms is not required to be
strictly associative, but it is required to be associative up to a
\emph{coherent} 2-isomorphism.  This means that for every composable
triple $\Theta$, $\Psi$ and $\Phi$ of 1-morphisms there is a specified
2-isomorphism $(\Theta\circ\Psi)\circ\Phi
\stackrel\sim\Longrightarrow\Theta\circ(\Psi\circ\Phi)$ and these
2-isomorphisms have to satisfy the so-called pentagon coherency
condition which ensures that although $\Theta\circ\Psi\circ\Phi$ is
ambiguous, it can be taken to mean either $(\Theta\circ\Psi)\circ\Phi$
or $\Theta\circ(\Psi\circ\Phi)$ without confusion.  The up-shot of
this is that parentheses are unnecessary in the notation.


Each object $X$ also comes with an identity 1-morphism $\id_X$, but
again, in general, one does not have equality of ${\id_Y}\circ \Phi$,
$\Phi$ and $\Phi\circ\id_X$, but rather the identity 1-morphisms come
with coherent 2-isomorphisms ${\id_Y}\circ\Phi\stackrel\sim
\Rightarrow\Phi$, and $\Phi\circ\id_X\stackrel\sim\Rightarrow\Phi$.
Again this means that in practice the identities can be neglected in
the notation: so although we could denote the identity 1-morphism
with, say, a dotted line, we choose not to. This is illustrated below.
\[\pic{SD6a}\equiv\pic{SD6b}\]


A \emph{strict} 2-category is one in which the coherency
2-isomorphisms for associativity and identities are themselves all
identities.  So the 2-category $\gCat$ of categories, functors and
natural transformations is a strict 2-category.


The last piece of structure that a 2-category has is the
\emph{horizontal composition} of 2-morphisms.  If $\Phi,\Phi'\colon
X\to Y$ and $\Psi,\Psi'\colon Y\to Z$ are 1-morphisms, and
$\alpha\colon \Phi\Rightarrow\Phi'$ and $\beta\colon
\Psi\Rightarrow\Psi'$ are 2-morphisms, then $\beta\circh \alpha\colon
\Psi\circ\Phi\Rightarrow\Psi'\circ\Phi'$ is defined, and is notated as
below.
\[\pic{AD7a}\equiv\pic{AD7b}\quad\pic{SD7a}\equiv\pic{SD7b}\]
The horizontal and vertical composition are required to obey the
\emph{interchange law}.
\[(\beta'\circv\beta)\circh(\alpha'\circv\alpha)=
(\beta'\circh\alpha')\circv(\beta\circh\alpha).\]
This means that the following diagrams are unambiguous.
\[\pic{AD8}\qquad\pic{SD8}\]
It also means that 2-morphisms can be `slid past' each other in the
following sense.
\[\pic{SD9a}=\pic{SD9b}\]

From now on, string diagrams will be drawn without the grey borders, and labels will be omitted if they are clear from the context.


 \subsection{The 2-category $\gVar$.}
The 2-category $\gVar$, of spaces and integral kernels, is defined as
follows.  The objects are spaces, as defined in the introduction, and the hom-category $\bHom_\gVar(X,Y)$ from a space $X$
to a space $Y$ is the derived category $\D(X\times Y)$, which is to be
thought of as the category of integral kernels from $X$ to $Y$.
Explicitly, this means that the 1-morphisms in $\gVar$ from $X$ to $Y$
are objects of $\D(X\times Y)$ and the 2-morphisms from $\Phi$ to
$\Phi'$ are morphisms in $\Hom_{\D(X\times Y)}(\Phi,\Phi')$, with
vertical composition of 2-morphisms just being usual composition in
the derived category.

Composition of 1-morphisms in $\gVar$ is defined using the convolution
of integral kernels: if $\Phi\in \D(X\times Y)$ and $\Psi\in
\D(Y\times Z)$ are 1-morphisms then define the convolution
$\Psi\circ\Phi\in\D(X\times Z)$ by
\[ \Phi\circ \Psi := \pi^{}_{XZ, *} ( \pi_{YZ}^* \Psi \otimes
\pi_{XY}^* \Phi), \] where $\pi_{XZ}$, $\pi_{XY}$ and $\pi_{YZ}$ are the
projections from $X\times Y\times Z$ to the appropriate factors.  The
horizontal composition of 2-morphisms is similarly defined.  Finally,
the identity 1-morphism $\id_X\colon X\ra X$ is given by
$\cO_{\Delta}\in \D(X\times X)$, the structure sheaf of the diagonal
in $X\times X$.


The above 2-category is really what is at the heart of the study of
integral transforms, and it is entirely analogous to $\gAlg$, the
2-category of algebras described above.  For example, the Hochschild
cohomology groups of a space $X$ arise as the second homotopy groups
of the 2-category $\gVar$, at $X$:
\begin{align*}
  \HH^\blob(X) &:=\Ext_{X\times X}^\blob(\cO_\Delta,\cO_\Delta)
  \cong \Hom^\blob_{D(X\times X)}(\cO_\Delta,\cO_\Delta)\\
  &=:\THom_{\gVar}(\Id_X,\Id_X).
\end{align*}

There is a 2-functor from $\gVar$ to $\gCat$ which encodes integral
transforms: this 2-functor sends each space $X$ to its derived
category $\D(X)$, sends each kernel $\Phi\colon X\to Y$ to the
corresponding integral transform $\overline \Phi\colon D(X)\to D(Y)$,
and sends each map of kernels to the appropriate natural
transformation.  Many of the statements about integral transforms have
better formulations in the language of the 2-category $\gVar$.
